% ===== 原书 p.19 =====（§2.3 末尾在 p.19 上半页，已由 c02 处理；本块自 §2.4 起）
\subsection{奇维流形上的指标}

设 $D:\Gamma E\to\Gamma F$ 是紧奇维流形 $X$ 上的椭圆偏微分算子。前面已经指出，椭圆偏微分算子的象征是正齐次映射。用局部坐标写出，对一切实数 $\lambda>0$ 都有 $\sigma(x,\lambda\xi)=\lambda^m\sigma(x,\xi)$，其中 $m$ 是算子的阶。然而象征在局部是多项式，它满足更强的条件：对任意 $\lambda$（包括负值）都有 $\sigma(x,\lambda\xi)=\lambda^m\sigma(x,\xi)$。特别地，$\sigma(x,-\xi)=\pm\sigma(x,\xi)$，其中符号在整个流形上为常数。

于是记 $\alpha:TX\to TX$ 为逐纤维对径映射（丛对合），它由每根纤维上乘以 $-1$ 给出：$\alpha:(x,\xi)\mapsto(x,-\xi)$。考虑诱导映射 $\alpha^*:K(TX)\to K(TX)$。把象征类 $\sigma(D)\in K(TX)$ 用复形
\[
0\to\pi^*E\xrightarrow{\ \sigma(x,\xi)\ }\pi^*F\to 0
\]
表示，则元素 $\alpha^*\sigma(P)$ 由复形
\[
0\to\pi^*E\xrightarrow{\ \sigma(x,-\xi)\ }\pi^*F\to 0
\]
表示。由此得到 $\alpha^*\sigma(P)=\sigma(P)$：若 $\sigma(x,-\xi)=+\sigma(x,\xi)$，这立即成立；若 $\sigma(x,-\xi)=-\sigma(x,\xi)$，则把两个复形沿单位圆旋转半圈，得到它们之间的同伦。

\transnote{「旋转半圈」之所以有效，是因为在复线上 $e^{i\theta}$ 当 $\theta$ 从 $0$ 走到 $\pi$ 时始终可逆，把 $1$ 连续变到 $-1$；若改在实线上做同样的事，途中必然经过 $0$。两个复形因此同伦等价，它们在 $K(TX)$ 中的类相同。}

\begin{theorem}
紧奇维流形上任意椭圆偏微分算子的指标为零。
\end{theorem}

\begin{proof}[证明 1]
回想 Chern 特征给出从有理 K 理论到有理上同调的自然同构，有下面的交换图：
\[
\begin{tikzcd}
K^*(TX)\otimes\Q \arrow[r,"\alpha^*"] \arrow[d,"\ch"{right},"\cong"{left}]
  & K^*(TX)\otimes\Q \arrow[d,"\ch"{right},"\cong"{left}] \\
H^*_c(TX;\Q) \arrow[r,"\alpha^*"']
  & H^*_c(TX;\Q)
\end{tikzcd}
\]
由上同调的 Thom 同构定理，$H^*_c(TX;\Q)$ 是由 Thom 类 $u$ 生成的 $H^*(X;\Q)$ 模；$u$ 是唯一在每个 $TX$ 纤维上限制为选定定向类的上同调类。又因 $X$ 是奇维的，对径映射 $\alpha:TX\to TX$ 反转每根纤维的定向，从而 $\alpha^*u=-u$。

\transnote{图中竖直的 $\ch$ 是 Chern 特征给出的有理同构，两个水平箭头都写作 $\alpha^*$，但分别作用在 K 理论与上同调上。这张图交换，说的正是 $\ch$ 的自然性；证明 1 不过是沿这个方块走一圈。}

% ===== 原书 p.20 =====
由此得到，对每个 $x\in H^*_c(TX;\Q)$ 都有 $\alpha^*x=-x$，因此对应的 K 理论映射也由 $\alpha^*x=-x$ 给出，其中 $x\in K^*(TX)\otimes\Q$。既然 $\alpha^*\sigma(D)=\sigma(D)$，$\sigma(D)$ 在 $K^*(TX)\otimes\Q$ 中的像为零，从而 $\sigma(D)$ 是 $K(TX)$ 中的有限阶元素。拓扑指标是从 $K(TX)$ 到整数的同态，因此 $\Ind D=\tind\sigma(D)=0$。
\end{proof}

\transnote{同态 $\tind:K(TX)\to\Z$ 的像无挠，有限阶元素只能映到 $0$，这是证明 1 的收尾。注意这一步只用到 $\sigma(D)$ 在有理化之后为零，并没有说 $\sigma(D)$ 本身为零。}

\begin{proof}[证明 2]
由指标定理的上同调形式 (A)，有
\[
\Ind D=(-1)^n\big(\ch(\sigma(d))\smile\IndexClass(X)\big)[TX],
\]
其中 $n$ 是 $X$ 的维数。再作用对径映射 $\alpha:TX\to TX$，得到
\[
\begin{aligned}
\Ind D&=(-1)^n\big(\ch(\alpha^*\sigma(D))\smile\IndexClass(X)\big)_{\alpha_*[TX]}\\
&=(-1)^n\big(\ch(\sigma(D))\smile\IndexClass(X)\big)(-1)^n[TX]\\
&=-\Ind D.
\end{aligned}
\]
注意指标类 $\IndexClass(X)=\Td(TX\otimes_\R\C)\in H^*(X,\Q)$ 在 $\alpha^*$ 的作用下不变：复化映射 $\alpha:TX\otimes_\R\C\to TX\otimes_\R\C$ 沿单位圆旋转半圈即同伦于恒等映射。于是 $\Ind D=0$。
\end{proof}

注意，这两个证明的关键事实都是：在任意奇维流形上，对径映射反转 $TX$ 的定向。还要注意，证明 1 说明 $K(TX)$ 的每个元素未必都是某个椭圆偏微分算子的象征。特别地，在圆周上 $K(TS^1)\cong\tilde K(S^2)\cong\Z$，而每个椭圆偏微分算子的象征都是零。

\subsection{de Rham 复形}

现在再来看（复化后的）de Rham 复形 $D^\bullet$：
\[
0\to\Gamma\Lambda^0(TX\otimes\C)\xrightarrow{d}\Gamma\Lambda^1(TX\otimes\C)
\xrightarrow{d}\cdots\xrightarrow{d}\Gamma\Lambda^n(TX\otimes\C)\to 0.
\]
由 de Rham 定理，该复形的欧拉示性数为
\[
\chi(D^\bullet)=\sum_i(-1)^i\dim H^i(D^\bullet)=\sum_i(-1)^i\dim H^i_{dR}(X)=\chi(X).
\]
先考虑 $X$ 为偶数维 $n=2l$ 的情形。为计算拓扑指标，注意由 Euler 类的初等性质有
\[
c_n(TX\otimes\C)=e\big((TX\otimes\C)_\R\big)=(-1)^le(TX\oplus TX)=(-1)^le(X)^2,
\]
其中必须引入符号 $(-1)^l$，因为 $(TX\otimes\C)_\R$ 与 $TX\oplus TX$ 的定向不同。\footnote{设 $\{x_1,\ldots,x_n\}$ 是实向量空间 $V$ 的有序基。取 $\{x_1,ix_1,\ldots,x_n,ix_n\}$ 作为 $(V\otimes\C)_\R$ 的基，取 $\{x_1,\ldots,x_n,x'_1,\ldots,x'_n\}$ 作为 $V\oplus V$ 的基。若 $n=2l$，这两个基之间的置换符号为 $(-1)^l$。}

现在用一个稍作推广的分裂原理，允许（在偶维情形）假设 $TX\otimes\C$ 分解为线丛直和，其形式为
% ===== 原书 p.21 =====
\[
TX\otimes\C=(L_1\oplus\bar L_1)\oplus\cdots\oplus(L_l\oplus\bar L_l).
\]
记 $c_1(L_i)=x_i$，注意 $c_1(\bar L_i)=-x_i$，于是可以算出：
\[
c_n(TX\otimes\C)=\prod_i x_i\smile(-x_i)=(-1)^l\prod_i x_i^2,
\]
\[
\IndexClass(X)=\Td(TX\otimes\C)=\prod_i\frac{x_i}{1-e^{-x_i}}\frac{-x_i}{1-e^{x_i}},
\]
\[
\begin{aligned}
\ch\Big(\sum_{i=0}^n(-1)^i\Lambda^i(TX\otimes\C)\Big)
&=\ch\Big(\prod_i(1-L_i)(1+L_i)\Big)\\
&=\prod_i(1-e^{x_i})(1-e^{-x_i})\\
&=c_n(TX\otimes\C)\smile\Td(TX\otimes\C)^{-1}\\
&=(-1)^le(X)^2\smile\Td(TX\otimes\C)^{-1}.
\end{aligned}
\]
注意下式
\[
\ch\Big(\sum_{i=0}^n(-1)^i\Lambda^i(TX\otimes\C)\Big)\smile e(X)^{-1}
=(-1)^le(X)\smile\Td(TX\otimes\C)^{-1}
\]
在 $e(X)=0$ 时仍良定义，于是可以套用 Atiyah--Singer 指标定理的上同调形式 (C)，得到
\[
\chi(X)=\chi(D^\bullet)=(-1)^{n(n+1)}(-1)^le(X)[X]=e(X)[X].
\]

\transnote{原书此处印作 $(-1)^{n(n+1)}$，而定理 C 中的因子是 $(-1)^{n(n+1)/2}$。取 $n=2l$ 代入可见两者相差 $(-1)^l$，只有取半指数时才有 $\chi(X)=e(X)[X]$。译文照排原文，引用时请以定理 C 为准。}

在奇维情形，Euler 类 $e(X)$ 为零；由上面关于奇维流形上微分算子指标的讨论，$\chi(D^\bullet)=0$ 也成立。于是对任意流形 $X$，de Rham 复形的 Atiyah--Singer 指标定理给出
\[
\xi(X)=e(X)[X].
\]

\subsection{Dolbeault 复形}

现在考察 de Rham 复形的复类比。设 $X$ 是复 $n$ 维流形，记 $T_\C X$ 为 $X$ 的复切丛（即关于 $X$ 上复结构的切丛），有典范同构
\[
TX\otimes\C\cong T_\C X\oplus\overline{T_\C X}.
\]
由此得到分解
\[
\Lambda^i(TX\otimes\C)\cong\sum_{p+q=i}\Lambda^p(\overline{T_\C X})\otimes\Lambda^q(T_\C X)
=\sum_{p+q=i}\Lambda^{p,q}(T_\C X),
\]
其中定义
\[
\Lambda^{p,q}(T_\C X)=\Lambda^p(\overline{T_\C X})\otimes\Lambda^q(T_\C X).
\]
由此可见，$X$ 上复微分形式的空间 $\Gamma\Lambda^i(TX\otimes\C)$ 有分解
\[
\Gamma\Lambda^i(TX\otimes\C)=\sum_{p+q=i}\Gamma\Lambda^{p,q}(T_\C X),
\]
其中 $\Gamma\Lambda^{p,q}(T_\C X)$ 称为 $(p,q)$ 型光滑微分形式的空间。

取 $X$ 上的局部复坐标 $z_1,\ldots,z_n$，其中 $z_j=x_j+iy_j$，$(p,q)$ 型微分形式为
\[
\omega=\sum_{\substack{j_1<\cdots<j_p\\ k_1<\cdots<k_q}}a_{J,K}\,
dz_{j_1}\cdots dz_{j_p}\,d\bar z_{k_1}\cdots d\bar z_{k_q},
\]
这里 $J=(j_1,\ldots,j_p)$ 与 $K=(k_1,\ldots,k_q)$ 是多重指标，$\bar z_j=x_j-iy_j$ 表示复共轭。de Rham 算子 $d:\Gamma\Lambda^i(TX\otimes\C)\to\Gamma\Lambda^{i+1}(TX\otimes\C)$ 于是分解为
% ===== 原书 p.22 =====
$d=\partial+\bar\partial$。这里 $\partial:\Gamma\Lambda^{p,q}(T_\C X)\to\Gamma\Lambda^{p+1,q}(T_\C X)$ 与 $\bar\partial:\Gamma\Lambda^{p,q}(T_\C X)\to\Gamma\Lambda^{p,q+1}(T_\C X)$ 由以下的局部公式定义：
\[
\partial:\sum_{J,K}a_{J,K}dz_Jd\bar z_K\mapsto
\sum_{j,J,K}\frac{\partial}{\partial z_j}a_{J,K}\,dz_jdz_Jd\bar z_K,
\]
\[
\bar\partial:\sum_{J,K}a_{J,K}dz_Jd\bar z_K\mapsto
\sum_{k,J,K}\frac{\partial}{\partial\bar z_k}a_{J,K}\,d\bar z_kdz_Jd\bar z_K,
\]
其中
\[
\frac{\partial}{\partial z_j}=\frac12\Big(\frac{\partial}{\partial x_j}-i\frac{\partial}{\partial y_j}\Big),
\qquad
\frac{\partial}{\partial\bar z_j}=\frac12\Big(\frac{\partial}{\partial x_j}+i\frac{\partial}{\partial y_j}\Big).
\]
把 de Rham 复形限制到 $(0,q)$ 型微分形式，得到 Dolbeault 复形，它是如下椭圆复形：
\[
0\to\Gamma\Lambda^{0,0}(T_\C X)\xrightarrow{\bar\partial}\Gamma\Lambda^{0,1}(T_\C X)
\xrightarrow{\bar\partial}\cdots\xrightarrow{\bar\partial}\Gamma\Lambda^{0,n}(T_\C X)\to 0.
\]

为计算 Dolbeault 复形的指标，注意 $X$ 是复流形，Euler 类为 $e(X)=e(TX)=c_n(T_\C X)$。因此可以套用指标定理的形式 (C)。按定义有 $\Lambda^{0,q}(T_\C X)=\Lambda^q(T_\C X)$。

\transnote{这一步是定义层面的恒等：$\Lambda^{0,q}(T_\C X)=\Lambda^0(\overline{T_\C X})\otimes\Lambda^q(T_\C X)$，而 $\Lambda^0$ 是平凡线丛，故只剩 $\Lambda^q(T_\C X)$。不要把 $\Lambda^{0,q}$ 读成外代数里某个新的直和项。}

用分裂原理，不妨设 $T_\C X=L_1\oplus\cdots\oplus L_n$，$x_i=c_1(L_i)$，算出：
\[
\ch\Big(\sum_{i=0}^n\Lambda^i(T_\C X)\Big)=\ch\Big(\prod_i(1-L_i)\Big)=\prod_i(1-e^{x_i}),
\]
\[
e(X)=c_n(T_\C X)=\prod_i x_i,
\]
\[
\IndexClass(X)=\Td(TX\otimes\C)=\Td(T_\C X)\smile\Td(\overline{T_\C X})
=\Td(T_\C X)\smile\prod_i\frac{-x_i}{1-e^{x_i}}.
\]
把这些表达式代入指标定理，得到
\[
\begin{aligned}
\chi(\bar\partial)&=(-1)^{2n(2n+1)/2}
\Big(\ch\Big(\sum_{i=0}^n(-1)^i\Lambda^i(T_\C X)\Big)\smile e(X)^{-1}\smile\Td(TX\otimes\C)\Big)[X]\\
&=(-1)^{n(2n+1)}
\Big(\prod_i\frac{1-e^{x_i}}{x_i}\smile\Td(T_\C X)\smile\prod_i\frac{-x_i}{1-e^{x_i}}\Big)[X]\\
&=(-1)^n\big((-1)^n\Td(T_\C X)\big)[X]=\Td(T_\C X)[X].
\end{aligned}
\]
推广这一构造，可以得到著名的 Hirzebruch--Riemann--Roch 定理。完整的讨论见 [11, \S4]。

\begin{example}
把球面 $S^2$ 看作复 $1$ 维流形（即 Riemann 球），Dolbeault 复形约化为 $S^2$ 上的 Dolbeault 算子
\[
\bar\partial:\Gamma\Lambda^0(T_\C S^2)\to\Gamma\Lambda^1(T_\C S^2),
\]
它由 $\bar\partial f=(\partial/\partial\bar z)\,d\bar z$ 给出。由上面的讨论有
\[
\Ind\bar\partial=\Td(T_\C S^2)[X].
\]
现在回想 Todd 类由
\[
\Td(E)=1+\frac12c_1(E)+(\text{高阶项})
\]
给出，
% ===== 原书 p.23 =====
代入后得到
\[
\begin{aligned}
\Ind\bar\partial&=\Td(T_\C S^2)[X]=\big((1+\tfrac12c_1(T_\C S^2)\big)[X]\\
&=\tfrac12\big(c_1(T_\C S^2)\big)[X]=\tfrac12e(X)[X],
\end{aligned}
\]
而上一步刚证明 $e(X)[X]=\chi(X)$，于是
\[
\Ind\bar\partial=\frac12\chi(S^2)=1.
\]
更一般地，设 $X$ 是亏格 $g$ 的 Riemann 面，则 $\chi(X)=2-2g$，从而
\[
\Ind\bar\partial_X=1-g.
\]
\end{example}
